\section{Proof of Joint Selection Along a Schedule}
\label{sec:phil-joint-selection-proof}

\noindent\emph{Proof of \Cref{cor:phil-joint-selection}.}
The saddle estimate \cref{eq:phil-selection-saddle-spectrum} gives
\(\alpha(\mathrm T)\ge7/5\) throughout the stated temperature range.
The proof of \Cref{conj:phil-contrast-formation} permits any window
constant greater than \(1/\alpha(\mathrm T)\), independently of the
initial mean. Thus the common choice \(2\) is valid. Write
\[
 \tau_\varepsilon=2\varepsilon^{-1}\log(1/\varepsilon),
 \qquad W_\varepsilon=[\tau_\varepsilon,2\tau_\varepsilon],
 \qquad 0<\varepsilon<e^{-1}.
\]
For a perfect matching \(M\), define
\[
 \nu_{\mathrm T}^M
 =\mathcal N(0,\mathrm T\mathsf A_M^{-1})\otimes\delta_M,
 \qquad
 \mu_{\mathrm T}=\mu_{\mathrm T,1/2}.
\]
Here \(\mu_{\mathrm T,p}\) is the frozen-row law used in the preceding
proof, with its temperature now displayed explicitly. The inverse
exists because \(\mathsf A_M\succeq\mathrm TI\). Let
\(M_{\mathrm T}(\mathsf m_0)\) be the matching selected by the sign of
\(\mathsf j_{\mathrm T}(\mathsf m_0)\), choosing either matching on its
null zero set.

First record the full one-time consequence of the fixed-temperature
proof. Conditional on assignment \(M\), the law
\(\mu_{\mathrm T,p}\) has the Gaussian factor of \(\nu_{\mathrm T}^M\),
independently of \(p\). The normalized sheet weights are at most the
matching weight one. For
\(\zeta_i=|p_i-p_i^M|\), therefore,
\[
 \mu_{\mathrm T,p}(M)\ge\pi_M\ge1-\sum_i\zeta_i,
 \qquad
 \operatorname{TV}(\mu_{\mathrm T,p},\nu_{\mathrm T}^M)
 =1-\mu_{\mathrm T,p}(M)\le\sum_i\zeta_i.
\]
The row convergence and
\cref{eq:phil-selection-global-tracking} imply
\[
 \sup_{t\in W_\varepsilon}
 \operatorname{TV}
 (q_t^{\mathrm T,\varepsilon,\mathsf m_0},
       \nu_{\mathrm T}^{M_{\mathrm T}(\mathsf m_0)})
 \longrightarrow0
\]
at fixed temperature outside the imprint's null zero set.
Moreover \cref{eq:phil-selection-tracking} itself gives convergence
of the learned relative entropy to zero, uniformly on this window.
For frozen rows, the entropy identity and contraction give
\[
 \KL(q_t^{\mathrm T,0,\mathsf m_0}\Vert\mu_{\mathrm T})
 \le e^{-\lambda(\mathrm T,1/2)t}
       \KL(q_0^{\mathrm T,0,\mathsf m_0}\Vert\mu_{\mathrm T}).
\]
The initial entropy is finite by the Gaussian density bounds, and
\(\lambda(\mathrm T,1/2)>0\). Frozen entropy and total variation
therefore tend to zero on the window for every fixed mean.

Use the explicit entropy majorants as follows. Let
\[
 G_{\mathrm T}=\sup_{p\in[0,1]^4}\|\mathfrak g(p)\|_1,\qquad
 \ell_{\mathrm T}(\varepsilon)
 =\frac{\mathrm c_{\mathrm T}}
 {\mathrm K_{\mathrm T}+8\log2+32\log(1/\varepsilon)},\qquad
 \lambda_u(\mathrm T)
 =\frac{\mathrm c_{\mathrm T}}{\mathrm K_{\mathrm T}+8\log2}.
\]
The first constant is finite. The row lower bound and
\cref{eq:phil-selection-entropy-contraction} make
\(\ell_{\mathrm T}(\varepsilon)\) a rate lower bound throughout
\(t\le2\tau_\varepsilon\), and \(\lambda_u(\mathrm T)\) a rate lower
bound at uniform rows. The common initial entropy equals
\[
 H_0(\mathrm T,\mathsf m_0)
 =\log\mathsf Z(1/2)-6+
 \frac{\operatorname{tr}\overline{\mathsf A}
              +\mathsf m_0^\top\overline{\mathsf A}\mathsf m_0}
      {2\mathrm T},
 \qquad
 \overline{\mathsf A}=\frac1{16}\sum_\sigma\mathsf A_\sigma.
\]
Indeed the log density ratio is
\(\log\mathsf Z(1/2)-|\mathbf x-\mathsf m_0|^2/2+
\mathbf x^\top\mathsf A_\sigma\mathbf x/(2\mathrm T)\);
taking its initial expectation proves the formula.
Thus \(H_0\) is continuous and nonnegative. Uniformly on
\(W_\varepsilon\), the entropy tracking estimate gives
\[
 \begin{aligned}
 \KL(q_t^{\mathrm T,\varepsilon,\mathsf m_0}
               \Vert\mu_{\mathrm T,p(t)})
 &\le B_L(\mathrm T,\varepsilon,\mathsf m_0)\\
 &:=\left[
 e^{-\ell_{\mathrm T}(\varepsilon)\tau_\varepsilon/2}\sqrt{H_0}
 +\frac{\varepsilon G_{\mathrm T}}
             {\sqrt2\ell_{\mathrm T}(\varepsilon)}
 \right]^2,\\
 \KL(q_t^{\mathrm T,0,\mathsf m_0}\Vert\mu_{\mathrm T})
 &\le B_F(\mathrm T,\varepsilon,\mathsf m_0)
 :=e^{-\lambda_u(\mathrm T)\tau_\varepsilon}H_0.
 \end{aligned}
\]
These majorants are continuous in positive mobility and mean and tend
to zero at fixed temperature and mean as \(\varepsilon\downarrow0\).

Let \(\mathcal E_{\mathrm T}(\varepsilon,\mathsf m_0)\) be the minimum
of one and the maximum of the following five nonnegative quantities:
\[
 \begin{aligned}
 &\sup_{t\in W_\varepsilon}
          \max_i|p_i(t)-p_i^{M_{\mathrm T}(\mathsf m_0)}|,\\
 &\sup_{t\in W_\varepsilon}
       \operatorname{TV}
       (q_t^{\mathrm T,\varepsilon,\mathsf m_0},
                   \nu_{\mathrm T}^{M_{\mathrm T}(\mathsf m_0)}),\\
 &B_L(\mathrm T,\varepsilon,\mathsf m_0),\\
 &\sup_{t\in W_\varepsilon}
       \operatorname{TV}
          (q_t^{\mathrm T,0,\mathsf m_0},\mu_{\mathrm T}),\\
 &B_F(\mathrm T,\varepsilon,\mathsf m_0).
 \end{aligned}
\]
The preceding estimates prove
\(\mathcal E_{\mathrm T}(\varepsilon,\mathsf m_0)\to0\)
for almost every fixed mean at each fixed temperature.

We will take a supremum over mobility, so measurability must be checked.
Fix a temperature and restrict the positive mobility and the mean to
compact sets, with the mobility bounded away from zero. The corresponding
horizons are bounded, and
\cref{eq:phil-selection-row-bounds} bounds every row away from zero and
one. On this row set the posterior redraw probabilities are uniformly
Lipschitz in the rows, independently of the configuration.
Maximally couple the initial laws and then use common Brownian motions
and proposal clocks until the first different redraw. The law difference
is bounded by the initial total variation plus a constant times the
integral of row differences. Subtracting the learning equations bounds
the row difference by integrals of row and law differences and by the
mobility difference. Gr\"onwall's inequality proves total-variation
continuity in the initial mean and positive mobility on each bounded
time interval. Time continuity follows from
\cref{eq:phil-selection-forward}.

Parametrizing the window by \(t=u\tau_\varepsilon\),
\(1\le u\le2\), proves continuity in positive mobility of each row or
total-variation window supremum. The explicit entropy majorants are
continuous by their formulas. The matching as a function of the mean is
Borel because its imprint is analytic. Thus the error is measurable
in the mean, and its supremum over an interval of positive mobilities
equals its supremum over rational mobilities. No uniform continuity at
zero mobility or separate entropy-continuity assertion is needed.

Fix the sequence \(\mathrm T_n\downarrow0\) in the stated temperature
interval and \(0<\eta_n<1\) with \(\eta_n\to0\). For the prescribed
exponent \(a>7/2\), put \(h_n=\mathrm T_n^a\). The following argument
in fact permits arbitrary positive caps \(h_n\).
On the mean space \(\R^{12}\), introduce the auxiliary probability
measure \(\rho=\mathcal N(0,I_{12})\). Its density is strictly positive;
it introduces no assumption that the physical mean is random.
Set
\[
 F_n(r,\mathsf m_0)
 =\sup_{0<\varepsilon\le r}
       \mathcal E_{\mathrm T_n}(\varepsilon,\mathsf m_0).
\]
The measurable functions \(F_n\) are bounded by one. At fixed \(n\),
they tend to zero for \(\rho\)-almost every mean as \(r\downarrow0\).
Dominated convergence yields
\[
 \int F_n(r,\mathsf m_0)d\rho(\mathsf m_0)\longrightarrow0.
\]
Choose \(d_n=2^{-k_n}\) recursively, where \(k_n\) is the smallest
positive integer satisfying
\[
 d_n\le\min(e^{-1},2^{-n},h_n,d_{n-1}/2),
 \qquad
 \int F_n(d_n,\mathsf m_0)d\rho(\mathsf m_0)
       \le\eta_n2^{-n},
\]
with the \(d_{n-1}/2\) condition omitted for \(n=1\). The preceding
limit proves existence. This defines deterministic thresholds
independent of the mean; it does not give numerical values or an explicit
power law.

Let \(B_n=\{\mathsf m_0:F_n(d_n,\mathsf m_0)>\eta_n\}\).
Markov's inequality and countable subadditivity give
\[
 \rho(B_n)\le2^{-n},
 \qquad
 \rho\left(\bigcup_{n\ge N}B_n\right)\le2^{1-N}.
\]
The decreasing intersection of these tail unions consequently has
measure zero. Add the countable union of the imprint zero sets. The
resulting exceptional set is Lebesgue-null because \(\rho\) and Lebesgue
measure have the same null sets. Outside it, for every sufficiently large
\(n\), all \(0<\varepsilon\le d_n\) satisfy
\(F_n(d_n,\mathsf m_0)\le\eta_n\). Since \(\eta_n<1\), clipping has
no effect on this inequality: the row and total-variation errors and both entropy majorants are at
most \(\eta_n\). The majorants bound the actual entropies, giving
all five stated conclusions.

The frozen matching probability is, exactly,
\[
 \mu_{\mathrm T}(M)
 =\frac1{2+8\mathsf r_E+4\mathsf r_N+2\mathsf r_C}
 =\frac1{2D_0(\mathrm T)}
 =\frac12+O(\mathrm T).
\]
Indeed \(\mathsf r_E,\mathsf r_N=O(\mathrm T)\) and
\(\mathsf r_C=O(\mathrm T^{3/2})\), as proved from the sheet
determinants. Pushforward cannot increase total variation. If
\(q^{S}\) denotes the assignment marginal, the two established
total-variation estimates therefore imply, uniformly on \(W_\varepsilon\),
\[
 \left|
 \operatorname{TV}
 (q_t^{\mathrm T_n,\varepsilon,\mathsf m_0,S},
                  q_t^{\mathrm T_n,0,\mathsf m_0,S})
 -\left(1-\frac1{2D_0(\mathrm T_n)}\right)
 \right|\le2\eta_n.
\]
Here the comparison distance is exactly
\(\operatorname{TV}(\delta_M,\mu_{\mathrm T_n}^S)
=1-\mu_{\mathrm T_n}(M)\), for either matching.
The assignment-law contrast thus tends to \(1/2\), even if the selected
matching changes with \(n\).

Finally, defining \(\mathrm T(\varepsilon)=\mathrm T_n\) on
\(d_{n+1}<\varepsilon\le d_n\) gives a single stepwise path with both
parameters tending to zero. The proved bounds hold for all sufficiently
small mobilities on this path. Arbitrary mobility caps are admissible;
for example \(h_n=\mathrm T_n^n\) forces
\(\varepsilon_n/\mathrm T_n^a\to0\) for each fixed \(a>0\) whenever
\(\varepsilon_n\le d_n\). This imposes an upper bound and does not certify
a prescribed equality between temperature and mobility.

Nothing in this argument makes the imprint sign eventually constant as
temperature vanishes. Its fixed-temperature nonvanishing and null zero
set give no quantitative lower bound at a fixed mean in that limit.
Accordingly the matching is allowed to depend on \(n\), and the result
concerns complete one-time laws, not complete histories or conditional
first-formation events. \(\square\)




